\documentclass[11pt]{article}
\usepackage[fleqn]{amsmath}
\usepackage{amssymb, amsthm}
\usepackage[utf8]{inputenc}
\usepackage[margin=1in]{geometry}
\usepackage{graphicx}
\usepackage{enumitem}
\theoremstyle{plain}
\newtheorem{proposition}{Proposition}
\newcommand{\fit}[1]{\resizebox{\ifdim\width>\linewidth\linewidth\else\width\fi}{!}{#1}}
\setlength{\parindent}{0pt}
\begin{document}
\begin{proposition}[clear-pivot-cross]
Suppose $a \text{ s.t. } \operatorname{is\text{-}euclidean\text{-}ring}(a)$. Let $p \in \mathbb{N}, q \in \mathbb{N}, m \in \operatorname{mat}(\operatorname{succ}(p), \operatorname{succ}(q), \operatorname{carr}(a))$. If $\exists \mathit{i0}.\; \exists \mathit{j0}.\; ((\mathit{i0} \in \operatorname{interval}(1, \operatorname{succ}(p))) \wedge ((\mathit{j0} \in \operatorname{interval}(1, \operatorname{succ}(q))) \wedge \neg (\operatorname{entry}(m, \mathit{i0}, \mathit{j0}) = \operatorname{zero}(a))))$, then 
\begin{equation*}
\begin{array}{@{}l@{}} \exists \mathit{c2}. \\ \quad \operatorname{mat\text{-}equiv}(a, \operatorname{succ}(p), \operatorname{succ}(q), m, \mathit{c2}) \wedge \\ \quad \mathit{c2} \\ \quad =\; \operatorname{border}(a, \operatorname{entry}(\mathit{c2}, 1, 1), \operatorname{submat}(\mathit{c2}, p, q), p, q) \wedge \\ \quad \neg (\operatorname{entry}(\mathit{c2}, 1, 1) = \operatorname{zero}(a)) \end{array}
\end{equation*}

\end{proposition}

\begin{proof}
\begin{itemize}[leftmargin=2.8em, itemsep=4pt, parsep=1pt]
  \item[\textbf{8}] By $\operatorname{euclidean\text{-}ring\text{-}is\text{-}integral\text{-}domain}$:
\begin{equation*}
\operatorname{is\text{-}integral\text{-}domain}(a)
\end{equation*}

  \item[\textbf{9}] By $\operatorname{integral\text{-}domain\text{-}is\text{-}commutative\text{-}ring}$:
\begin{equation*}
\operatorname{is\text{-}commutative\text{-}ring}(a)
\end{equation*}

  \item[\textbf{10}] By $\operatorname{commutative\text{-}ring\text{-}is\text{-}ring}$:
\begin{equation*}
\operatorname{is\text{-}ring}(a)
\end{equation*}

  \item[\textbf{11--14}] \emph{Establish in-carrier / typing side-conditions:}
\begin{equation*}
\begin{array}{@{}l@{}}
(1 \in \operatorname{interval}(1, (p + 1))) \\
(1 \in \operatorname{interval}(1, (q + 1))) \\
((p + 1) \in \mathbb{N}) \\
((q + 1) \in \mathbb{N})
\end{array}
\end{equation*}

  \item[\textbf{15}] By $\operatorname{class\text{-}min\text{-}pivot}$:
\begin{equation*}
\begin{array}{@{}l@{}} \exists b. \\ \quad \operatorname{mat\text{-}equiv}(a, (p + 1), (q + 1), m, b) \wedge \\ \quad \neg (\operatorname{entry}(b, 1, 1) = \operatorname{zero}(a)) \wedge \\ \quad \forall c. \\ \quad \quad \operatorname{mat\text{-}equiv}(a, (p + 1), (q + 1), m, c) \Rightarrow  \\ \quad \quad \quad \forall i. \forall j. \\ \quad \quad \quad \quad (i \in \operatorname{interval}(1, (p + 1))) \Rightarrow  \\ \quad \quad \quad \quad \quad (j \in \operatorname{interval}(1, (q + 1))) \Rightarrow  \\ \quad \quad \quad \quad \quad \quad \neg (\operatorname{entry}(c, i, j) = \operatorname{zero}(a)) \Rightarrow  \\ \quad \quad \quad \quad \quad \quad \quad \operatorname{gauge}(a)(\operatorname{entry}(b, 1, 1)) \\ \quad \quad \quad \quad \quad \quad \quad \leq\; \operatorname{gauge}(a)(\operatorname{entry}(c, i, j)) \end{array}
\end{equation*}

  \item[\textbf{16--19}] By a cut on the auxiliary claim.
  \item[\textbf{20--26}] By $\operatorname{mat\text{-}equiv\text{-}trans}$:
\begin{equation*}
\operatorname{mat\text{-}equiv}(a, (p + 1), (q + 1), m, c)
\end{equation*}

  \item[\textbf{27--28}] Instantiate the hypothesis:
\begin{equation*}
\begin{array}{@{}l@{}} \forall j. \\ \quad (\mathit{ii} \in \operatorname{interval}(1, (p + 1))) \Rightarrow  \\ \quad \quad (j \in \operatorname{interval}(1, (q + 1))) \Rightarrow  \\ \quad \quad \quad \neg (\operatorname{entry}(c, \mathit{ii}, j) = \operatorname{zero}(a)) \Rightarrow  \\ \quad \quad \quad \quad \operatorname{gauge}(a)(\operatorname{entry}(\operatorname{b\_506}, 1, 1)) \\ \quad \quad \quad \quad \leq\; \operatorname{gauge}(a)(\operatorname{entry}(c, \mathit{ii}, j)) \end{array}
\end{equation*}

  \item[\textbf{29--30}] Instantiate the hypothesis:
\begin{equation*}
\begin{array}{@{}l@{}} \operatorname{gauge}(a)(\operatorname{entry}(\operatorname{b\_506}, 1, 1)) \\ \leq\; \operatorname{gauge}(a)(\operatorname{entry}(c, \mathit{ii}, \mathit{jj})) \end{array}
\end{equation*}
  \emph{(holds by assumption)}
  \item[\textbf{31}] \emph{Establish in-carrier / typing side-conditions:}
\begin{equation*}
\begin{array}{@{}l@{}}
(\operatorname{b\_506} \in \operatorname{mat}((p + 1), (q + 1), \operatorname{carr}(a)))
\end{array}
\end{equation*}

  \item[\textbf{32}] By $\operatorname{clear\text{-}first\text{-}row}$:
\begin{equation*}
\begin{array}{@{}l@{}} \exists \operatorname{q\_509}. \\ \quad \operatorname{mat\text{-}equiv}(a, (p + 1), (q + 1), \operatorname{b\_506}, \operatorname{q\_509}) \wedge \\ \quad \operatorname{entry}(\operatorname{q\_509}, 1, 1) \\ \quad =\; \operatorname{entry}(\operatorname{b\_506}, 1, 1) \wedge \\ \quad \forall j \in \operatorname{interval}(1, (q + 1)). \\ \quad \quad \neg (j = 1) \Rightarrow  \\ \quad \quad \quad (\operatorname{entry}(\operatorname{q\_509}, 1, j) = \operatorname{zero}(a)) \end{array}
\end{equation*}

  \item[\textbf{33--36}] By $\operatorname{classmin\text{-}transport}$:
\begin{equation*}
\begin{array}{@{}l@{}} \forall c. \\ \quad \operatorname{mat\text{-}equiv}(a, (p + 1), (q + 1), \operatorname{q\_509\_511}, c) \Rightarrow  \\ \quad \quad \forall \mathit{ii}. \forall \mathit{jj}. \\ \quad \quad \quad (\mathit{ii} \in \operatorname{interval}(1, (p + 1))) \Rightarrow  \\ \quad \quad \quad \quad (\mathit{jj} \in \operatorname{interval}(1, (q + 1))) \Rightarrow  \\ \quad \quad \quad \quad \quad \neg (\operatorname{entry}(c, \mathit{ii}, \mathit{jj}) = \operatorname{zero}(a)) \Rightarrow  \\ \quad \quad \quad \quad \quad \quad \operatorname{gauge}(a)(\operatorname{entry}(\operatorname{q\_509\_511}, 1, 1)) \\ \quad \quad \quad \quad \quad \quad \leq\; \operatorname{gauge}(a)(\operatorname{entry}(c, \mathit{ii}, \mathit{jj})) \end{array}
\end{equation*}

  \item[\textbf{37}] \emph{Establish in-carrier / typing side-conditions:}
\begin{equation*}
\begin{array}{@{}l@{}}
(\operatorname{q\_509\_511} \in \operatorname{mat}((p + 1), (q + 1), \operatorname{carr}(a)))
\end{array}
\end{equation*}

  \item[\textbf{38}] By a cut on the auxiliary claim.
  \item[\textbf{39--40}] Substitute the established equation.  \emph{(holds by assumption)}
  \item[\textbf{41}] By $\operatorname{clear\text{-}first\text{-}col}$:
\begin{equation*}
\begin{array}{@{}l@{}} \exists \operatorname{q\_518}. \\ \quad \operatorname{mat\text{-}equiv}(a, (p + 1), (q + 1), \operatorname{q\_509\_511}, \operatorname{q\_518}) \wedge \\ \quad \operatorname{entry}(\operatorname{q\_518}, 1, 1) \\ \quad =\; \operatorname{entry}(\operatorname{q\_509\_511}, 1, 1) \wedge \\ \quad \forall i \in \operatorname{interval}(1, (p + 1)). \\ \quad \quad \neg (i = 1) \Rightarrow  \\ \quad \quad \quad (\operatorname{entry}(\operatorname{q\_518}, i, 1) = \operatorname{zero}(a)) \wedge \\ \quad \forall c \in \operatorname{interval}(1, (q + 1)). \\ \quad \quad \operatorname{entry}(\operatorname{q\_518}, 1, c) \\ \quad \quad =\; \operatorname{entry}(\operatorname{q\_509\_511}, 1, c) \end{array}
\end{equation*}

  \item[\textbf{42--46}] By $\operatorname{mat\text{-}equiv\text{-}trans}$:
\begin{equation*}
\operatorname{mat\text{-}equiv}(a, (p + 1), (q + 1), m, \operatorname{q\_509\_511})
\end{equation*}

  \item[\textbf{47}] By $\operatorname{mat\text{-}equiv\text{-}trans}$:
\begin{equation*}
\operatorname{mat\text{-}equiv}(a, (p + 1), (q + 1), m, \operatorname{q\_518\_520})
\end{equation*}

  \item[\textbf{48}] \emph{Establish in-carrier / typing side-conditions:}
\begin{equation*}
\begin{array}{@{}l@{}}
(\operatorname{q\_518\_520} \in \operatorname{mat}((p + 1), (q + 1), \operatorname{carr}(a)))
\end{array}
\end{equation*}

  \item[\textbf{49}] By a cut on the auxiliary claim.
  \item[\textbf{50}] Substitute the established equation.
  \item[\textbf{51--52}] Substitute the established equation.  \emph{(holds by assumption)}
  \item[\textbf{53}] By a cut on the auxiliary claim.
  \item[\textbf{54--57}] Instantiate the hypothesis:
\begin{equation*}
(\operatorname{entry}(\operatorname{q\_509\_511}, 1, \mathit{jr}) = \operatorname{zero}(a))
\end{equation*}

  \item[\textbf{58--59}] Substitute the established equation.  \emph{(holds by assumption)}
  \item[\textbf{60}] By $\operatorname{bordered\text{-}eq\text{-}border}$:
\begin{equation*}
\begin{array}{@{}l@{}} \operatorname{q\_518\_520} \\ =\; \operatorname{border}(a, \operatorname{entry}(\operatorname{q\_518\_520}, 1, 1), \operatorname{submat}(\operatorname{q\_518\_520}, p, q), p, q) \end{array}
\end{equation*}

  \item[\textbf{61--66}] Take $\mathit{c2} :=$ $\operatorname{q\_518\_520}$.  \emph{(holds by assumption)}
\end{itemize}
\end{proof}
\end{document}
